\documentclass[10pt,a4paper]{article}
\usepackage[utf8]{inputenc}
\usepackage[T1]{fontenc}
\usepackage[brazilian]{babel}
\usepackage{geometry}
\usepackage{helvet}
\usepackage[table]{xcolor}
\usepackage{enumitem}
\usepackage{array}
\usepackage{tabularx}
\renewcommand{\familydefault}{\sfdefault}
\geometry{top=1.8cm,bottom=1.8cm,left=2cm,right=2cm}
\setlength{\parindent}{0pt}
\setlength{\parskip}{0.4em}
\definecolor{azul}{HTML}{1F4E79}
\definecolor{azulclaro}{HTML}{EAF2F8}
\definecolor{cinza}{HTML}{F4F6F7}
\newcommand{\termo}[2]{\vspace{0.2em}\noindent\colorbox{azulclaro}{\parbox{\dimexpr\linewidth-2\fboxsep}{\textbf{#1} --- #2}}\par\vspace{0.2em}}
\newcommand{\secao}[1]{\vspace{0.75em}\noindent\textcolor{azul}{\Large\textbf{#1}}\par\vspace{0.18em}\hrule\vspace{0.32em}}
\pagestyle{plain}

\begin{document}

\begin{center}
{\LARGE\textbf{Apostila de Estudo}}\\[0.15cm]
{\Large\textcolor{azul}{Métodos da Ciência Moderna}}\\[0.2cm]
\textit{Questão 2 da prova 2026.1 --- Galileu, Bacon, Descartes, empirismo, racionalismo e mecanicismo.}
\end{center}

\secao{O que a questão realmente pede}

Não basta associar um autor a uma palavra. A resposta precisa mostrar quatro relações: \textbf{quem é o autor}, \textbf{qual é seu método}, \textbf{como ele produz conhecimento} e \textbf{qual corrente representa}. O mecanicismo entra como a visão de natureza que torna importante observar, medir e explicar os fenômenos por leis.

\secao{O contexto comum: a natureza como máquina}

\termo{Mecanicismo}{Entende a natureza como uma máquina. Matéria e energia são seus constituintes básicos, e a matemática é a ferramenta para descrever seus fenômenos. Nessa visão, os acontecimentos materiais possuem relações causais e podem ser explicados ou previstos por cálculo.}

Essa concepção ajuda a entender a ciência moderna: se a natureza tem ordem e relações, o pesquisador pode procurar essas relações pela razão, pela observação, pelo experimento e pela matemática.

\secao{As três abordagens}

\termo{Descartes --- racionalismo --- dedução}{A razão é o principal caminho para o conhecimento. O pesquisador parte de princípios, postulados ou um problema geral, organiza o raciocínio de modo lógico e chega a conclusões particulares. O método dedutivo busca demonstrar e justificar pela coerência, consistência e não-contradição.}

\termo{Bacon --- empirismo --- indução experimental}{O conhecimento deve ser construído a partir de experimentos, registro repetido de dados e análise. Parte-se de casos particulares e, após verificar padrões, chega-se a generalizações. Para Bacon, o experimento é o objeto de análise e verificação.}

\termo{Galileu --- empirismo --- indução observacional e quantitativa}{A investigação começa na observação direta do fenômeno. A partir dela, extraem-se elementos para análise e experimentação. Galileu une observação sistemática, precisão, instrumentos, experimento e linguagem matemática.}

\secao{A diferença decisiva: Bacon e Galileu}

Ambos se relacionam ao empirismo e ao método indutivo. Portanto, ambos caminham dos casos particulares para generalizações. A diferença está no \textbf{início do processo}:

\begin{center}
\fcolorbox{azul}{cinza}{\parbox{0.9\linewidth}{\centering\textbf{Bacon:} experimento $\longrightarrow$ dados $\longrightarrow$ análise $\longrightarrow$ generalização\\[0.1cm]\textbf{Galileu:} observação direta $\longrightarrow$ análise $\longrightarrow$ experimentação e quantificação}}
\end{center}

\newpage
\secao{Comparação para usar na prova}

\renewcommand{\arraystretch}{1.08}
\begin{tabularx}{\linewidth}{>{\bfseries}p{1.9cm}>{\bfseries}p{2.6cm}p{3.2cm}X}
\hline
\rowcolor{azulclaro}Autor & Corrente/método & Ponto de partida & Como produz conhecimento \\
\hline
Descartes & Racionalismo; dedução & Razão, princípios e problema geral & Encadeia deduções lógicas até soluções particulares. \\
Bacon & Empirismo; indução & Experimento e dados repetidos & Analisa e verifica os dados para formular generalizações. \\
Galileu & Empirismo; indução & Observação direta do fenômeno & Analisa, experimenta, mede e descreve quantitativamente. \\
\hline
\end{tabularx}

\secao{Como as abordagens se relacionam}

Elas não precisam ser tratadas como inimigas. O material mostra que Isaac Newton combinou deduções matemáticas com induções extraídas de observações e experimentos. Assim, uniu o método empírico associado a Bacon ao método racional associado a Descartes. A lição é que a ciência moderna se fortalece pela articulação entre razão, experiência e matemática.

\secao{Resposta de 4 pontos: estrutura segura}

\begin{enumerate}[leftmargin=0.7cm,itemsep=0.32em]
\item Abra definindo o mecanicismo como a visão de uma natureza ordenada, causal e matematicamente descritível.
\item Explique Descartes: racionalismo e método dedutivo, da razão e dos princípios gerais às conclusões particulares.
\item Explique Bacon: empirismo e indução, com experimentos, dados repetidos, análise e generalização.
\item Explique Galileu: observação direta, seguida de análise, experimento, medição e quantificação.
\item Feche mostrando que as três abordagens contribuíram para a ciência moderna e que Newton combinou razão e experiência.
\end{enumerate}

\textbf{Resposta-modelo:} O mecanicismo entende a natureza como uma máquina, formada por relações causais que podem ser descritas pela matemática. Nesse contexto, Descartes representa o racionalismo e o método dedutivo: a partir da razão e de princípios gerais, constrói demonstrações lógicas que conduzem a conclusões particulares. Bacon representa o empirismo e o método indutivo, baseado em experimentos, registro repetido de dados, análise e generalização. Galileu também se relaciona ao empirismo e à indução, mas inicia pela observação direta do fenômeno, seguindo para análise, experimentação e descrição quantitativa. Essas abordagens contribuíram para a ciência moderna ao valorizar razão, observação, experimento e matemática.

\secao{Variações de pergunta: o que adaptar}

\begin{tabularx}{\linewidth}{>{\bfseries}p{5cm}X}
\hline
\rowcolor{azulclaro}Se a pergunta pedir... & A ideia que não pode faltar \\
\hline
``Compare Bacon e Galileu'' & Ambos são indutivos e empiristas; Bacon começa pelo experimento, Galileu pela observação direta. \\
``Explique racionalismo'' & Razão e dedução: princípios gerais conduzem a conclusões particulares. \\
``Relacione mecanicismo e ciência moderna'' & A natureza é considerada ordenada, causal e descritível pela matemática. \\
``Como a ciência moderna se formou?'' & Ela articulou razão, experiência, observação, experimento e quantificação; Newton simboliza essa combinação. \\
\hline
\end{tabularx}

\secao{Teste de domínio}

Você domina a questão se consegue explicar, sem consultar: por que Bacon e Galileu não são idênticos; por que Descartes é racionalista; e como o mecanicismo conecta matemática e explicação da natureza.

\vfill
\small\textcolor{azul}{Atalho mental: Descartes raciocina e deduz; Bacon experimenta e generaliza; Galileu observa, mede e testa.}

\end{document}
